\documentclass[a4paper,11pt]{article}
\usepackage[utf8]{inputenc}
\usepackage[T1]{fontenc}
\usepackage[brazil]{babel}
\usepackage{longtable}
\usepackage{array}
\usepackage[margin=2.5cm]{geometry}

\begin{document}

\section*{Plano de Indexação 1}

\begin{longtable}{|l|p{7cm}|l|l|}
\hline
\textbf{Tabela} & \textbf{Justificativa} & \textbf{Variáveis} & \textbf{Tipo} \\
\hline
\endfirsthead
\hline
\textbf{Tabela} & \textbf{Justificativa} & \textbf{Variáveis} & \textbf{Tipo} \\
\hline
\endhead
\texttt{tb\_evento} & A maioria dos relatórios financeiros calcula a receita a partir do preço do evento. Operações de \texttt{MAX(vl\_preco)}, ordenação e comparação por preço aparecem em consultas de público de alto valor e análises de receita. & \texttt{vl\_preco} & B-Tree \\
\hline
\texttt{rl\_inscricao} & Praticamente todos os relatórios financeiros consideram apenas inscrições com pagamento \textit{CONCLUIDO}. O filtro por status é o mais recorrente em consultas de receita real, perda potencial e preferência de pagamento. & \texttt{cd\_status\_pagamento} & Hash \\
\hline
\end{longtable}

\bigskip
\section*{Plano de Indexação 2}

\begin{longtable}{|l|p{7cm}|l|l|}
\hline
\textbf{Tabela} & \textbf{Justificativa} & \textbf{Variáveis} & \textbf{Tipo} \\
\hline
\endfirsthead
\hline
\textbf{Tabela} & \textbf{Justificativa} & \textbf{Variáveis} & \textbf{Tipo} \\
\hline
\endhead
\texttt{tb\_evento} & A maioria dos relatórios financeiros calcula a receita a partir do preço do evento. Operações de \texttt{MAX(vl\_preco)}, ordenação e comparação por preço aparecem em consultas de público de alto valor e análises de receita. & \texttt{vl\_preco} & Hash \\
\hline
\texttt{rl\_inscricao} & Praticamente todos os relatórios financeiros consideram apenas inscrições com pagamento \textit{CONCLUIDO}. O filtro por status é o mais recorrente em consultas de receita real, perda potencial e preferência de pagamento. & \texttt{cd\_status\_pagamento} & Hash \\
\hline
\end{longtable}

\bigskip
\section*{Plano de Indexação 3}

\begin{longtable}{|l|p{7cm}|l|l|}
\hline
\textbf{Tabela} & \textbf{Justificativa} & \textbf{Variáveis} & \textbf{Tipo} \\
\hline
\endfirsthead
\hline
\textbf{Tabela} & \textbf{Justificativa} & \textbf{Variáveis} & \textbf{Tipo} \\
\hline
\endhead
\texttt{tb\_participante} & O CPF é o identificador único de pessoas no sistema e usado em buscas exatas (lookup por documento). Decidimos utilizar o index hash pois a variável CPF segue uma distribuição linear. & \texttt{nu\_cpf} & Hash \\
\hline
\texttt{tb\_revisor} & Mesma justificativa: localização pontual de revisores pelo CPF em validações e consultas de conflito de interesse, sempre via igualdade. & \texttt{nu\_cpf} & Hash \\
\hline
\end{longtable}

\end{document}
